\documentclass{article}

\usepackage{amsmath,amssymb,amsthm}

\newtheorem{definition}{Definition}

\newtheorem{theorem}{Theorem}

\newtheorem{lemma}{Lemma}

\newtheorem{proposition}{Proposition}

\newtheorem{corollary}{Corollary}

\newtheorem{example}{Example}

\newtheorem{remark}{Remark}

\newtheorem{problem}{Problem}

\begin{document}

Throughout this note, all variables range over the integers.

\begin{definition}

\label{def:even}

An integer $n$ is even if there exists $k\in\mathbb{Z}$ such that $n=2k$.

\end{definition}

\begin{definition}

\label{def:odd}

An integer $n$ is odd if there exists $k\in\mathbb{Z}$ such that $n=2k+1$.

\end{definition}

\newpage

\begin{theorem}

\label{thm:add}

If $a$ and $b$ are even, then $a+b$ is even.

\end{theorem}

Thus the set of even integers is closed under addition.

\begin{lemma}

\label{lem:negative}

If $a$ is even, then $-a$ is even.

\end{lemma}

\begin{proof}

By Definition 1, write $a=2m$. Then $-a=2(-m)$ and $-m\in\mathbb{Z}$.

\end{proof}

\newpage

\begin{proposition}

\label{prop:product}

If $a$ is even and $b\in\mathbb{Z}$, then $ab$ is even.

\end{proposition}

\begin{proof}

By Definition 1, $a=2m$ for some integer $m$. Hence $ab=2(mb)$, so the product is even.

\end{proof}

\begin{theorem}

\label{thm:odd-sum}

The sum of two odd integers is even.

\end{theorem}

\begin{proof}

By Definition 2, write $a=2m+1$ and $b=2n+1$. Then $a+b=2(m+n+1)$, which is even by Definition 1.

\end{proof}

\newpage

\begin{corollary}

\label{cor:finite}

Every finite sum of even integers is even.

\end{corollary}

\begin{proof}

Apply Theorem 1 inductively to the summands. The empty sum is $0=2\cdot 0$.

\end{proof}

\begin{example}

\label{ex:add}

As an instance of Theorem 1, $2+4=6$ is even.

\end{example}

\begin{remark}

\label{rem:multiples}

Even integers are precisely the integer multiples of $2$.

\end{remark}

\begin{problem}

\label{problem:odd-product}

Prove that the product of two odd integers is odd.

\end{problem}

\newpage

\begin{proof}[Proof of Theorem~\ref{thm:add}]

By Definition 1, there exist $m,n\in\mathbb{Z}$ with $a=2m$ and $b=2n$.

Thus $a+b=2(m+n)$. Since $m+n\in\mathbb{Z}$, Definition 1 gives the conclusion.

\end{proof}

\end{document}
